% LSPV-5 BINARYDOF_UNIFORM1 — 2026-10-02
\section*{Binary Degrees of Freedom and Uniform Worst-Case Solver Capacity}

\textbf{STATUS: CANDIDATE / OPEN.} This note proposes a framework and does not prove a lower bound or separate P from NP.

\subsection*{1. Source-conditioned irreducible binary freedom}
Distinguish raw semantic assignment freedom from computationally irreducible binary freedom. A provisional diagnostic is
\[
D^*(\Phi;M)=\min_{R\in\mathcal R_M(\Phi)}D(\Phi\mid R).
\]
The research target is an independently characterizable source invariant that lower-bounds irreducible freedom without defining hardness into the measure.

Existing characteristic-two observables include
\[
D_{\rm rank}=\operatorname{rank}_{\mathbb F_2}\Phi_{\rm bin},
\qquad
D_{\rm null}=\binom m2-D_{\rm rank},
\]
MATRIXCERT1 legal-slice innovations, and EXCESS/TURANEX1 source-sensitive deficits
\[
\epsilon(I)=d-\operatorname{rank}G_{E(I)}.
\]
An experimental ledger should track explicit freedom, source rank, excess, legal-slice innovations, replay-response classes, and the best known compressed representation.

\subsection*{2. Lipschitz form}
If an independently defined invariant satisfies
\[
D(\Phi_0)\ge B(n),\quad D(\Phi_{\rm final})=O(1),\quad
D_{t+1}\ge D_t-L(n),
\]
then
\[
T(n)\ge \frac{B(n)-O(1)}{L(n)}.
\]
A superpolynomial lower bound requires a superpolynomial ratio. Explicitly exponential trees or response sets do not suffice when they have polynomial-size generators.

\subsection*{3. Per-instance shortcut versus one legitimate solver}
For each individual input there may exist a short branch, certificate, restriction sequence, or source-specific representation that rapidly determines the answer:
\[
\forall x\;\exists R_x\quad \operatorname{Cost}(x,R_x)\le\operatorname{poly}(|x|).
\]
This does not imply
\[
\exists A\;\forall x\quad \operatorname{Cost}_A(x)\le\operatorname{poly}(|x|).
\]
The quantifier order is load-bearing.

A deterministic solver cannot select an answer-dependent shortcut merely because one exists. It must compute which shortcut to use, or maintain enough uniform structure to cover every input. Acquisition, indexing, routing, normalization, and verification of the shortcut must therefore be charged.

This resembles the distinction between an accepting branch of a nondeterministic computation and a deterministic algorithm that must locate or bypass the information supplied by that branch. An NDTM does not literally know the answer: NP acceptance is existential. The unresolved deterministic question is whether the relevant existential choice can be found or bypassed in polynomial time.

\subsection*{4. Circuit version}
For fixed input length n, a circuit C_n must compute the correct function on all 2^n inputs. An input-specific simplified circuit C_x, branch B_x, or restriction path is not an upper bound for the whole function unless the collection can be selected or combined within the allowed size.

Separate
\[
\text{path complexity}(x),\qquad
\text{selection/routing complexity}(x),\qquad
\text{global solver representation complexity}(n).
\]
A candidate accounting ledger is
\[
\operatorname{Cost}_{\rm solver}
=
\operatorname{Cost}_{\rm local}
+\operatorname{Cost}_{\rm acquire/select}
+\operatorname{Cost}_{\rm route}
+\operatorname{Cost}_{\rm shared},
\]
with reusable shared information charged only once.

The desired phenomenon is a family for which local shortcuts exist but no polynomial-size common selector, router, or shared representation covers all worst-case inputs.

\subsection*{5. Direct-product and replay connection}
A desired theorem could resemble
\[
D^*(A_1\otimes\cdots\otimes A_k)\ge\sum_iD^*(A_i).
\]
But class count or explicit branching is insufficient. One must rule out a succinct common evaluator mapping source plus challenge history directly to the correct response. STRATFLAT1 and LEGALFLAT1 remain mandatory hostile controls.

\subsection*{6. Proposed uniform-coverage experiment}
On one exact reconstruction-compatible SU/vMS source family:
\begin{enumerate}
\item Enumerate legal source instances, including the mandatory n=5 case.
\item Compute the cheapest known local shortcut or response representation for each instance.
\item Cluster instances by common shortcut/evaluator.
\item Compute the minimum number and total description size of common evaluators covering the source family.
\item Charge the selector mapping an input/source to its evaluator.
\item Compare with universal interpolation, Gaussian elimination, Skolem circuits, decision diagrams, CRT/arithmetic straight-line programs, and source/history quotients.
\item Repeat after legal restrictions and measure whether coverage complexity grows or collapses to a compact quotient.
\end{enumerate}

The central quantity is not the number of local branches but the complexity of one answer-independent mechanism covering all of them.

\subsection*{7. Hostile controls}
Parity/Tseitin with Gaussian elimination; Horn, dual-Horn and 2-SAT; bounded treewidth; independent products; padding/renaming/auxiliary variables; DEAD/COPY/DELAY/INTERP/RECOMP; memoization; static reread; batching; pointer jumping; semantic and algebraic oracle replacement; succinct circuits/branching programs; and nonuniform advice.

The circuit formulation must distinguish uniform algorithms from nonuniform circuit families. A lower bound against uniform selection alone is not an unrestricted circuit lower bound.

\subsection*{8. Candidate proof shape}
The strongest route would establish an independently defined D with
\[
\text{explicit NP-complete source}\Longrightarrow D(\Phi_n)\ge B(n),
\]
\[
\text{every legal deterministic step/shared representation}\Longrightarrow |\Delta D|\le L(n),
\]
and
\[
B(n)/L(n)=n^{\omega(1)}.
\]
For circuits one additionally needs a theorem forcing enough shared routing/representation capacity to realize the relevant binary distinctions.

No such theorem is established. The immediate task is to test whether BITRANK1, MATRIXCERT1, EXCESS/TURANEX1, and source-bound replay measure aspects of one source-conditioned irreducible binary-freedom object, or whether compact universal evaluators collapse it.

\subsection*{Strict conclusion}
\[
\boxed{\forall x\;\exists\text{ cheap shortcut for }x}
\quad\not\Rightarrow\quad
\boxed{\exists\text{ one cheap deterministic solver covering all }x}.
\]
This is a quantifier and acquisition distinction, not yet a lower bound. The burden is to prove that covering or selecting all local shortcuts requires superpolynomial resources without answer-dependent advice or importing nondeterministic choice into the deterministic model.

P versus NP remains OPEN.

% END_OF_NOTE::BINARYDOF_UNIFORM1
